\section{Incidence expansion}\label{sec:incidence}

Let $S$ be a finite set of distinct points in an affine plane, and let a
\emph{secant} be a line containing at least two of the points. Write
$\mathcal L$ for the secants and $t_i$ for the number through point $i$.
For centered real labels $v_i$, meaning $\sum_i v_i=0$, we have
\[
 \sum_{L\in\mathcal L}\left|\sum_{i\in S\cap L}v_i\right|^2
       =\sum_{i\in S}(t_i-1)|v_i|^2 .
\]
On expanding the squares, each diagonal term occurs $t_i$ times and
each ordered pair of distinct points occurs once. Centering gives the
identity. We need a version for points placed in a spectral metric space,
where one can measure how far an entire secant pulls a common center
in a given direction. This version will bound the mean squared length
of an incidence edge in the graphs used for permutation routing
(Section~\ref{sec:sort-action}).

\subsection{An inequality for secant lines}

Take $S\subset F^2$ with $|S|\ge2$. For each $i\in S$ choose a site
$x_i$ in a complete $\mathrm{CAT}(0)$ spectral space $X$ over $F$
and a mass $a_i>0$. Let $o$ be the minimizer of
$\sum_i a_i d(o,x_i)^2$ and set $r_i=d(o,x_i)$. Two points of $S$
may be placed at the same site.

Suppose the legs for a subset $C\subset S$ have widths $a_i r_i$.
Sending them from a common site $m$ to their destinations costs
$\sum_{i\in C}a_i r_i d(m,x_i)$. By first variation at $m=o$, the
supremum of the initial savings per unit distance, allowing the common
site to stay at $o$, is
\[
 K_C=\sup_{y\ne o}
       \left(\sum_{i\in C}a_i r_i\cos\angle_o(x_i,y)\right)_+ .
\]
Here $y\in X$, $\angle_o$ is the Alexandrov angle and the summand is
zero if $r_i=0$. For a one-point space set $K_C=0$.
Write $K_L=K_{S\cap L}$. On a real line this is the absolute value of
the secant sum above, with $v_i=a_i(x_i-o)$.

\begin{theorem}[Incidence inequality]\label{thm:incidence}
If $F$ is algebraically closed and $\omega_F<9/4$, then for every
such finite configuration of affine points, weighted sites and secants,
\begin{equation}\label{eq:inc-secants}
       \sum_{L\in\mathcal L}K_L^2
          \le \sum_{i\in S}(t_i-1)a_i^2r_i^2 .
\end{equation}
\end{theorem}

\begin{proof}
Write $i=(u_i,v_i)$ and project the points to the sphere positions
$z_i(s)=u_i s+v_i$. As $s$ varies, a pair collides exactly once on
$\mathbb P^1(F)$, with the projection $i\mapsto u_i$ used at infinity.
Let $D$ consist of the finite collision parameters together with
$\infty$. For each $\alpha\in D$ let $\mathcal P_\alpha$ be the
partition into fibers of its projection. The nonsingleton fibers are
precisely the intersections with secants, each occurring at the
parameter for its direction. Thus $i$ is a singleton at exactly
$|D|-t_i$ parameters, and $K_{\{i\}}=a_i r_i$. Adding these singleton
contributions to both sides of~\eqref{eq:inc-secants} gives the
equivalent estimate
\begin{equation}\label{eq:inc-fibers}
   \sum_{\alpha\in D}\sum_{C\in\mathcal P_\alpha} K_C^2
             \le (|D|-1)\sum_i a_i^2 r_i^2 .
\end{equation}

\emph{Choosing the external widths.}
Fix $s_0\notin D$. The positions $z_i(s_0)$ are distinct, so
Theorem~\ref{thm:sphere} with $\lambda_i=1/(2a_i)$ gives energies
$h_i=h_i(N)$ with
\begin{equation}\label{eq:inc-saturation}
 w_N(h_i)\longrightarrow a_i r_i,\qquad
 \frac{1}{N}H_{\mathbf x}\bigl(C(\mathbf h;\mathbf z(s_0))\bigr)
              \longrightarrow P:=\sum_i a_i r_i^2 .
\end{equation}
Here $w_N(h)=c\sqrt h/N$, $c=\pi/\sqrt3$. The cost of supplying the
external legs from $o$, in these limiting widths, is exactly $P$.
If $P=0$ all $K_C$ vanish. Otherwise take $N$ large enough that the
chosen tensor is nonzero. Its total energy $H_0=\sum_i h_i$ is
$O(N^2)$ and is an integer, since every generator in any nonzero
coefficient is paired.

\emph{A joint bound on the collision energies.}
Write $T(s)=C(\mathbf h;\mathbf z(s))$. In the chart at a finite
$\alpha$ use $s-\alpha$ as coordinate. At infinity use $u=1/s$ and
the positions $u_i+u v_i$. Homogeneity in the positions gives the
tensor $u^{-H_0}T(1/u)$ in this chart. In either chart the points of
each fiber $C$ have a common center and distinct linear offsets.
Proposition~\ref{prop:fusion} expands the coefficient of order $l$
as a sum over intermediate energies $e_C\ge0$ satisfying
$\sum_C e_C=H_0+l$.

Choose temporarily a linear functional $\ell$ with
$f(s)=\ell(T(s))$ nonzero. Let $l_\alpha$ be the order of its first
nonzero coefficient in the corresponding tensor chart. In particular,
$l_\infty=\operatorname{ord}_\infty f-H_0$. The ordinary orders of
$f$ on $\mathbb P^1$ sum to zero. Since there are no poles off $D$,
the energy totals at these orders satisfy
\begin{equation}\label{eq:inc-orders}
  \sum_{\alpha\in D}(H_0+l_\alpha)\le (|D|-1) H_0 .
\end{equation}
Each summand is nonnegative. Using the \emph{same} functional for
all directions is what gives one bound on their total intermediate
energy.

We now require the selected coefficients to retain the price
in~\eqref{eq:inc-saturation}. In every chart collect the fusion terms
with total intermediate energy at most $(|D|-1)H_0$. There are only
polynomially many in $N$. Each term is specified by one nonnegative
half-integer per fiber, and the number of fibers and directions is
fixed. Let $W$ be the span of all collected terms whose price is
less than $H_{\mathbf x}(T(s_0))-\sqrt N$. The ordinary-sum inequality
bounds the price of every nonzero tensor in $W$ by this threshold
plus $O(\log N)$. For large $N$, $T(s_0)\notin W$. Take $\ell$
vanishing on $W$ and nonzero on $T(s_0)$. At each $\alpha$ the
coefficient at its first nonzero order $l_\alpha$ has a fusion term
surviving $\ell$. By~\eqref{eq:inc-orders} such a term was collected,
so its price is at least $H_{\mathbf x}(T(s_0))-\sqrt N$. Select one
surviving term per direction.

Pass to a common subsequence along which the selected widths
$w_N(e_C)$ converge to limits $b_{\alpha,C}$. There are finitely
many widths, all bounded by~\eqref{eq:inc-orders}. Multiplication
of that inequality by $c^2/N^2$ gives
\[
 \sum_{\alpha\in D}\sum_{C\in\mathcal P_\alpha}b_{\alpha,C}^2
                  \le (|D|-1)\sum_i a_i^2 r_i^2 .
\]

\emph{The width needed by one fiber.}
Fix $\alpha$, a fiber $C$, and $y\ne o$. Supply the selected term
by preparing the fused tensor at $o$. Carry the intermediate leg
for $C$ a distance $t$ along $[o,y]$ to a site $m$, apply its fusion
map there and send the external legs from $m$ to their sites. Split
all other legs at $o$. By the transport bound~\eqref{eq:prices-transport}
and $\log\dim V_h\le c\sqrt h$, the price of the selected term is at
most the cost of this supply. After division by $N$, the $\sqrt N$
price loss vanishes. Passing to the limit using~\eqref{eq:inc-saturation} gives
\[
 P\le\sum_{i\notin C}a_i r_i^2+b_{\alpha,C}t+
                       \sum_{i\in C}a_i r_i d(m,x_i).
\]
At $t=0$ equality holds. The derivative from the right is therefore
nonnegative. First variation in a $\mathrm{CAT}(0)$ space gives
$b_{\alpha,C}\ge\sum_{i\in C}a_i r_i\cos\angle_o(x_i,y)$
\cite[Chapter~II.3]{BridsonHaefliger}. Taking the supremum of the
positive part yields $b_{\alpha,C}\ge K_C$, proving~\eqref{eq:inc-fibers}.
\end{proof}

\subsection{The sharp local link inequality}

Suppose $F$ contains $\mathbb F_q$. Take the $q^2$ points of the affine
plane $\mathbb F_q^2$ and retain the lines in $s$ parallel classes,
$1\le s\le q+1$. A point belongs to $s$ retained lines and each line
contains $q$ points. Place the vertices of the resulting bipartite
incidence graph in $X$. Let $e$ be the mean squared distance across a
uniform edge. Give each side half the total mass, uniformly distributed
over its vertices, and let $V$ be the minimum mean squared distance
of this distribution to a common center.

\begin{corollary}[Affine incidence links]\label{cor:inc-link}
Under the hypotheses of Theorem~\ref{thm:incidence},
\[
                   e\ge 2(1-s^{-1/2}) V .
\]
The coefficient is optimal in every target containing a nontrivial
geodesic segment. Omitting one parallel class gives $s=q$ and a
$q$-regular graph on $2q^2$ vertices.
\end{corollary}

\begin{proof}
Take $o$ to be the barycenter of the point sites alone. Write $r_i,r_L$
for the distances from $o$ of a point site $x_i$ and a line site
$y_L$. The mean squared distances on the two sides are
\[
             A=q^{-2}\sum_i r_i^2,\qquad
             B=(sq)^{-1}\sum_{L\ \mathrm{retained}} r_L^2 .
\]
For each incidence, triangle comparison bounds its squared distance
from below by $r_i^2+r_L^2-2r_i r_L\cos\angle_o(x_i,y_L)$.
A zero radius gives a zero cross term. By~\eqref{eq:inc-secants}
with all masses one, $\sum_L K_L^2\le q\sum_i r_i^2$ even for the
full set of lines. Thus the mean cross term over the $s q^2$
retained incidences is at most
\[
 \frac{2}{s q^2}\sum_{L\ \mathrm{retained}}r_L K_L
        \le\frac{2}{s q^2}
           \sqrt{q\sum_i r_i^2\sum_{L\ \mathrm{retained}}r_L^2}
        = 2\sqrt{AB/s}.
\]
Consequently
\[
      e\ge A+B-2\sqrt{AB/s}\ge (1-s^{-1/2})(A+B).
\]
Evaluation at $o$ gives $V\le(A+B)/2$, which proves the corollary.

To attain it, choose one retained parallel class and assign its $q$
lines real labels of mean zero and positive variance. Give each point
the label of its line. Give each line site of that class $\sqrt s$
times its label and every other line site zero. Both sides now have
mean zero and the same mean square $A$. A fraction $1/s$ of the edges
belongs to the chosen class, so the mean product of the two endpoint
labels is $A/\sqrt s$. Hence $V=A$ and
$e=2(1-s^{-1/2})A$. Translate and rescale the labels into any
nontrivial segment of $X$. Nearest-point projection onto the segment
ensures that the ambient minimum has the same value.
\end{proof}

The bound in terms of $A+B$ keeps the point barycenter fixed.
Section~\ref{sec:sort-action} uses that form for affine links, and
uses the secant inequality directly for projective links.
